\section{Module \texttt{toggle.c}}

Le module \texttt{toggle.c} couvre les composants binaires du wrapper : case à cocher et ligne composée avec interrupteur. Ces deux widgets reposent sur des contrôles GTK simples, mais le wrapper leur ajoute une configuration déclarative homogène.

\subsection{Fonction \texttt{create\_checkbox}}

\begin{lstlisting}[language=C]
GtkWidget *create_checkbox(const checkbox_config *config) {
    GtkWidget *checkbox = gtk_check_button_new_with_label(config && config->label ? config->label : "");

    if (config == NULL) {
        return checkbox;
    }

    gtk_check_button_set_active(GTK_CHECK_BUTTON(checkbox), config->active);
    gtk_check_button_set_inconsistent(GTK_CHECK_BUTTON(checkbox), config->inconsistent);
    apply_widget_style(checkbox, &config->style);

    if (config->on_toggled != NULL) {
        g_signal_connect(checkbox, "toggled", G_CALLBACK(config->on_toggled), config->user_data);
    }

    return checkbox;
}
\end{lstlisting}

\subsection{Fonction \texttt{create\_switch\_row}}

Cette seconde fonction construit un composant composite : un label à gauche, un \texttt{GtkSwitch} à droite, puis un éventuel callback sur la propriété \texttt{active}.

\begin{lstlisting}[language=C]
GtkWidget *create_switch_row(const switch_config *config, GtkWidget **out_switch) {
    GtkWidget *row = gtk_box_new(GTK_ORIENTATION_HORIZONTAL, 10);
    GtkWidget *label = gtk_label_new(config && config->label ? config->label : "");
    GtkWidget *sw = gtk_switch_new();

    gtk_widget_set_hexpand(label, TRUE);
    gtk_widget_set_halign(label, GTK_ALIGN_START);
    gtk_widget_set_halign(sw, GTK_ALIGN_END);
    gtk_box_append(GTK_BOX(row), label);
    gtk_box_append(GTK_BOX(row), sw);

    if (config != NULL) {
        gtk_switch_set_active(GTK_SWITCH(sw), config->active);
        gtk_switch_set_state(GTK_SWITCH(sw), config->state);
        apply_widget_style(row, &config->style);

        if (config->on_active_changed != NULL) {
            g_signal_connect(sw, "notify::active", G_CALLBACK(config->on_active_changed), config->user_data);
        }
    }

    if (out_switch != NULL) {
        *out_switch = sw;
    }

    return row;
}
\end{lstlisting}

\subsection{APIs GTK utilisées}

\begin{center}
\begin{tabular}{|l|p{8cm}|}
\hline
\textbf{API} & \textbf{Rôle} \\
\hline
\texttt{gtk\_check\_button\_new\_with\_label()} & Crée la case à cocher avec son libellé. \\
\hline
\texttt{gtk\_check\_button\_set\_inconsistent()} & Active l'état visuel intermédiaire. \\
\hline
\texttt{gtk\_switch\_new()} & Crée l'interrupteur binaire GTK4. \\
\hline
\texttt{gtk\_switch\_set\_active()} & Positionne l'état logique initial du switch. \\
\hline
\texttt{gtk\_box\_append()} & Assemble le label et l'interrupteur dans une même ligne. \\
\hline
\texttt{g\_signal\_connect()} & Relie les changements d'état aux callbacks utilisateur. \\
\hline
\end{tabular}
\end{center}

\newpage
